\documentclass[11pt,a4paper]{article}
\usepackage[utf8]{inputenc}
\usepackage[T1]{fontenc}
\usepackage[brazil]{babel}
\usepackage[margin=2.5cm]{geometry}
\usepackage{xcolor}
\usepackage{hyperref}

\definecolor{fjblue}{HTML}{0284C7}
\definecolor{fjdark}{HTML}{0F172A}

\hypersetup{
  colorlinks=true,
  linkcolor=fjblue,
  urlcolor=fjblue,
  pdftitle={Manual do Usuário — ForjaTikZ},
  pdfauthor={Jean Pimenta / Arkhé Code}
}

\title{\textbf{\color{fjdark}Manual do Usuário — ForjaTikZ}\\
\large Studio Visual de Diagramas TikZ e Gráficos PGFPlots (GNU GPL v3.0)}
\author{Jean Pimenta $\cdot$ Arkhé Code\\
\href{https://forjatex.com.br}{forjatex.com.br} $\cdot$ \href{https://github.com/jeanjpimenta/forjatex-suite}{GitHub}}
\date{Versão 1.1.0}

\begin{document}
\maketitle
\tableofcontents
\newpage

\section{Introdução}
O \textbf{ForjaTikZ} é um aplicativo visual complementar do \textbf{ForjaTeX App Suite} projetado para criar diagramas \textbf{TikZ} e gráficos matemáticos \textbf{PGFPlots} sem precisar memorizar toda a sintaxe do \LaTeX{}.

\section{Modo 1: Gráficos (PGFPlots)}
\begin{enumerate}
  \item \textbf{Tipo de Gráfico:} escolha entre \textbf{Funções $f(x)$}, \textbf{Pontos / CSV} (com suporte a barras de erro), \textbf{Barras (\texttt{ybar})} ou \textbf{Superfície 3D (\texttt{\textbackslash addplot3})}.
  \item \textbf{Configuração dos Eixos:} defina título, rótulos X e Y, limites (\texttt{xmin}, \texttt{xmax}, \texttt{ymin}, \texttt{ymax}), escala (\texttt{linear}, \texttt{semilogx}, \texttt{semilogy}, \texttt{loglog}) e grade.
  \item \textbf{Fórmulas Matemáticas:} escreva expressões em função de \texttt{x} (ex.: \texttt{0.2*x\^{}2}, \texttt{sin(x)}, \texttt{exp(-x\^{}2)}). Use os botões de atalho para inserir fórmulas prontas. Com a opção \textbf{Usar Radianos} ativada (\texttt{trig format plots=rad}), \texttt{sin(x)} funciona diretamente em radianos.
  \item \textbf{Dados e Barras:} no modo \textbf{Pontos / CSV}, insira um ponto por linha como \texttt{x, y} ou \texttt{x, y, erro}. No modo \textbf{Barras}, insira \texttt{Categoria=Valor} (ex.: \texttt{Jan=15}).
\end{enumerate}

\section{Modo 2: Diagramas (TikZ)}
\begin{itemize}
  \item \textbf{Formas (\texttt{shapes.geometric}):} adicione Círculos, Caixas, Losangos (decisão) e Elipses e arraste-os na grade com ajuste automático (\textit{snap}).
  \item \textbf{Conexões:} selecione um nó, clique em \textbf{Conectar Nó Selecionado} e clique no nó de destino. Edite rótulos e curvatura (\texttt{bend left} / \texttt{bend right}).
  \item \textbf{Modelos Prontos CTAN:} gere em 1 clique exemplos iniciais de \texttt{tikz-cd} (diagramas comutativos), \texttt{automata} (máquinas de estado) e \texttt{circuitikz} (circuitos).
\end{itemize}

\section{Exportação e Validação}
\begin{itemize}
  \item \textbf{Copiar Código:} copia o bloco \texttt{\textbackslash begin\{tikzpicture\}...\textbackslash end\{tikzpicture\}} para colar no seu documento no \textbf{ForjaTeX IDE}.
  \item \textbf{Exportar .tex:} salva o snippet em um arquivo \texttt{.tex}.
  \item \textbf{Testar LaTeX:} executa o \texttt{pdflatex} localmente (se instalado) para validar se o código compila sem erros.
\end{itemize}

\section{Glossário}
\begin{description}
  \item[TikZ / PGF] Linguagem vetorial nativa do \LaTeX{} para criação de diagramas e ilustrações técnicas.
  \item[PGFPlots] Pacote construído sobre TikZ especializado em gráficos científicos 2D e 3D.
  \item[\texttt{trig format plots=rad}] Opção do PGFPlots que faz funções trigonométricas aceitarem radianos em vez de graus.
\end{description}

\section{Histórico de Versões (Changelog)}
\begin{itemize}
  \item \textbf{v1.1.0:} Barra lateral corrigida, PGFPlots completo (Funções, Pontos/CSV + Barras de Erro, Barras, 3D, Log/Semilog), formas adicionais TikZ, modelos CTAN, seletor PT-BR/EN automático e envio de feedback para \texttt{forjatex.com.br}.
  \item \textbf{v1.0.0:} Lançamento inicial com canvas TikZ básico e gerador de funções PGFPlots.
\end{itemize}

\end{document}

